\chapter{Regeneration models: logistic, multiplicative, binomial}
\label{app:binomial}

This appendix derives the deterministic increment of Chapter~\ref{ch:env}
and the Stage~B multiplier of Chapter~\ref{ch:design} from a per-unit
regrowth process, states the renewable-resource models the thesis is
placed against, and reports the exact dynamic-programming values of every
policy in the main text under a third, Pérolat-faithful, regeneration
kernel.
Every number is printed by \texttt{python cpr\_xi.py} and pinned by
\texttt{tests/test\_cpr\_xi.py}.

\section{Per-unit regrowth and its mean field}
\label{app:meanfield}

In Commons Harvest \cite{perolat2017} an empty cell regrows with a
probability that rises with the number of apples nearby and is zero when
none are nearby.
The scalar version of that process on a pool of capacity \(K\) with
escapement (post-harvest stock) \(S\) is: each of the \(K-S\) empty units
regrows independently with probability \(\rho S/K\), so the increment is
\begin{equation}
  \begin{aligned}
  g_{\mathrm{bin}}(S)&\sim\operatorname{Binomial}\!\bigl(K-S,\;\rho S/K\bigr),\\
  \mathbb{E}\,g_{\mathrm{bin}}(S)&=\rho\,\frac{S\,(K-S)}{K},
  \qquad
  \operatorname{Var}g_{\mathrm{bin}}(S)=\rho\,\frac{S\,(K-S)}{K}\Bigl(1-\frac{\rho S}{K}\Bigr).
  \end{aligned}
  \label{eq:binomial}
\end{equation}
The mean is the logistic surplus-production increment of Schaefer
\cite{schaefer1954}, and \eqref{eq:growth} with \(\xi\equiv 1\) is its
rounding to the integers.
The deterministic game of Stage~A is therefore the mean-field limit of the
per-unit process, and Stage~B is a second-moment correction to it: the
multiplier \(\xi_t\in\{0.7,1.0,1.3\}\) preserves the mean of
\eqref{eq:binomial} exactly and adds a bounded, symmetric spread.
Table~\ref{tab:moments} compares the two spreads at the escapement levels
that decide the game.
The binomial spread is proportionally largest at low stock (a coefficient
of variation of \(0.53\) at \(S=4\), against \(0.27\) for \(\xi\)), which
is why the two models differ on the constant-pair cells below.

\begin{table}[ht]
\centering
\caption{Mean and standard deviation of the growth increment at
  escapement \(S\) under the executed \(\xi\) multiplier and under the
  binomial kernel \eqref{eq:binomial}.
  The exact mean is \(\rho S(K-S)/K\); the deterministic increment is its
  half-even rounding.
  \(S=4\) is the escapement under mutual take-2 from \(R_0\); \(S=5\) under
  hawk--dove; \(S=8\) after a symmetric restrained opening or at the joint
  optimum's escapement; \(S=14\) at the joint-optimum stock.}
\label{tab:moments}
\begin{tabular}{rrrrrrr}
\toprule
\(S\) & exact mean & deterministic & \(\xi\) mean & \(\xi\) s.d. & binomial mean & binomial s.d. \\
\midrule
4  & 3.24 & 3 & 3.00 & 0.82 & 3.24 & 1.72 \\
5  & 3.94 & 4 & 4.00 & 0.82 & 3.94 & 1.87 \\
8  & 5.76 & 6 & 5.67 & 1.25 & 5.76 & 2.17 \\
14 & 8.19 & 8 & 8.33 & 2.05 & 8.19 & 2.37 \\
20 & 9.00 & 9 & 9.00 & 2.45 & 9.00 & 2.22 \\
\bottomrule
\end{tabular}
\end{table}

\section{The renewable-resource models}
\label{app:reed}

The discrete-time harvesting model of Reed \cite{reed1979}, in the
notation of Clark \cite{clark1990}, is
\begin{equation}
  X_{t+1}=Z_t\,G(S_t),
  \qquad
  S_t=X_t-h_t,
  \qquad
  Z_t\ \text{i.i.d.},\ \mathbb{E}Z_t=1,
  \label{eq:reed}
\end{equation}
where \(X_t\) is the stock before harvest, \(h_t\) the harvest, \(S_t\)
the escapement, \(G\) a concave stock--recruitment function and \(Z_t\) a
multiplicative recruitment shock.
Reed's result is that the policy maximising expected discounted net
revenue is a constant-escapement rule, \(h_t=\max(0,X_t-S^{*})\), and that
under a self-sustaining condition (the escapement is reachable from every
realised stock) the stochastic optimum \(S^{*}\) coincides with the
deterministic one.
Sethi et al.\ \cite{sethi2005} separate three sources of uncertainty,
growth (\(Z_t\)), measurement (the manager observes \(X_t\) with error) and
implementation (the realised harvest differs from the intended one), and
show that the constant-escapement result survives growth uncertainty alone
but not large measurement or implementation uncertainty.

The game of Chapter~\ref{ch:env} is \eqref{eq:reed} with two harvesters,
\(h_t=a_{1,t}+a_{2,t}\) capped at 3 per player, the Schaefer recruitment
\(G(S)=S+\rho S(K-S)/K\) rounded to the integers, and the shock applied to
the surplus rather than to the whole recruitment:
\begin{equation}
  R_{t+1}=S_t+\xi_t\,g(S_t)
  \qquad\text{instead of}\qquad
  X_{t+1}=Z_t\bigl(S_t+g(S_t)\bigr).
  \label{eq:surplus-shock}
\end{equation}
The difference is deliberate.
Under \eqref{eq:reed} a small \(Z_t\) can take the stock below the
escapement, so collapse can be environmental; under
\eqref{eq:surplus-shock}, since \(\xi_t\ge 0\), the stock never falls below
the escapement and collapse is always a consequence of the players'
requests, which is what makes survival a behavioural statistic.
In Sethi et al.'s taxonomy the thesis studies growth uncertainty only:
the stock is reported to both players exactly every round, and requests
are filled exactly unless the scarcity rule fires.
This is the regime in which Reed's constant-escapement structure survives,
and the dynamic programme confirms it in the two-player game: the joint
optimum of \eqref{eq:bellman-joint} is the constant-escapement policy
``harvest to \(S=8\)'' (one round of \((0,0)\) from \(R_0=8\), then
\((3,3)\) at the fixed point \(R=14\)), and every safe policy of
Table~\ref{tab:noise-dp} is a feedback rule of the same form, constrained
by the per-player cap of 3, which binds below the maximum sustainable
yield of \(\rho K/4=9\).
Reed's model is thus not a form citation for Stage~B; it is the model, with
the shock moved from recruitment to surplus and the harvest split between
two players who do not coordinate.

\section{Dynamic-programming values under the three kernels}
\label{app:dp3}

Table~\ref{tab:dp3} reports every policy of Table~\ref{tab:noise-dp}
under the deterministic increment, the executed \(\xi\) multiplier, and
the binomial kernel \eqref{eq:binomial}.
Under the binomial kernel the constant-pair cells of Table~\ref{tab:chicken}
are no longer invariant: mutual take-2 survives with probability
\(0.17\) and hawk--dove dies with probability \(0.19\), so the opening
incentives of Stage~A would not carry over and the two arms would not be
comparable.
That is the reason the executed Stage~B uses the bounded multiplier
rather than the process it approximates.
The stock-conditioned feedback rules keep their values under all three
kernels (the symmetric rule ``3 if \(R\ge 14\), else 0'' earns
\(102.7\) per player with survival \(1.00\) under both stochastic kernels),
which is the constant-escapement result of Section~\ref{app:reed} in the
two-player game.
The best response to a partner that has learned the feedback rule is
\(107.0\), \(105.5\) and \(104.2\) under the three kernels, against
\(72\) for a hawk that stays on take-2; that gap is the ceiling for what a
shaper could gain by exploiting a partner it has taught to restrain, and
the shaper's return in Chapter~\ref{ch:results} is read against it.

\begin{table}[ht]
\centering
\small
\caption{Expected return per player and survival probability under the
  three regeneration kernels.
  ``Hawk'' is the frozen always-2 partner of Test~A.
  The \(\xi\) column is the executed Stage~B; the binomial column is the
  Pérolat-faithful reference of \eqref{eq:binomial}.
  Exact backward induction, rounded to one decimal; \((1,1)\) under the
  binomial kernel survives with probability \(0.999\).}
\label{tab:dp3}
\resizebox{\textwidth}{!}{%
\begin{tabular}{lcccccc}
\toprule
 & \multicolumn{2}{c}{Deterministic} & \multicolumn{2}{c}{\(\xi\in\{0.7,1,1.3\}\)} & \multicolumn{2}{c}{Binomial} \\
\cmidrule(lr){2-3}\cmidrule(lr){4-5}\cmidrule(lr){6-7}
Policy pair & Return & Survival & Return & Survival & Return & Survival \\
\midrule
\((1,1)\) throughout & 36 / 36 & 1.00 & 36.0 / 36.0 & 1.00 & 36.0 / 36.0 & 1.00 \\
\((2,2)\) throughout & 7 / 7 & 0.00 & 7.4 / 7.4 & 0.00 & 18.1 / 18.1 & 0.17 \\
\((2,1)\) hawk versus dove & 72 / 36 & 1.00 & 72.0 / 36.0 & 1.00 & 60.6 / 30.3 & 0.81 \\
\midrule
\((1,1)\) once, then \((2,2)\) & 71 / 71 & 1.00 & 59.3 / 59.3 & 0.80 & 51.6 / 51.6 & 0.68 \\
Hawk versus ``1 once, then 2'' & 72 / 71 & 1.00 & 35.7 / 34.7 & 0.40 & 35.3 / 34.3 & 0.42 \\
Hawk versus ``0 once, then 2'' & 72 / 70 & 1.00 & 60.3 / 58.3 & 0.80 & 52.6 / 50.6 & 0.68 \\
\((0,0)\) once, then \((3,3)\) & 105 / 105 & 1.00 & 38.6 / 38.6 & 0.25 & 46.3 / 46.3 & 0.36 \\
\midrule
Hawk versus ``1 if \(R<12\), else 2'' & 72 / 69 & 1.00 & 72.0 / 68.8 & 1.00 & 60.6 / 56.9 & 0.81 \\
Both ``2 if \(R\ge 9\), else 1'' & 71 / 71 & 1.00 & 70.8 / 70.8 & 1.00 & 69.7 / 69.7 & 0.99 \\
Both ``3 if \(R\ge 14\), else 0'' & 105 / 105 & 1.00 & 102.7 / 102.7 & 1.00 & 102.7 / 102.7 & 1.00 \\
\midrule
Best response to hawk & 105 & 1.00 & 102.5 & 1.00 & 101.0 & 1.00 \\
Best response to ``1 if \(R<12\), else 2'' & 107 & 1.00 & 105.5 & 1.00 & 104.2 & 1.00 \\
Joint optimum \(W^{*}\) & 210 & 1.00 & 207.1 & 1.00 & 206.3 & 1.00 \\
\bottomrule
\end{tabular}%
}
\end{table}
